\chapter{Notation Reference}
\label{app:math}

This appendix is a notation reference, not a tutorial. The two mathematical bridges at the start of the book develop these objects in context: Bridge~I introduces representation, vectors, matrices, norms, and geometry; Bridge~II introduces probability, loss, empirical risk, and gradients. The main chapters then build on that vocabulary. This appendix collects the resulting symbols and conventions in one place so that a reader who has worked through the bridges can look up a forgotten symbol without rereading the development. Entries point to the chapter or bridge where the object is introduced.

\section{Typographic Conventions}

The book uses the following conventions consistently. Where an entry states ``the bridges'' or a chapter label, that is where the convention is established and used.

\begin{itemize}[leftmargin=*]
\item \textbf{Scalars} are lowercase italic: $n$, $d$, $p$, $\lambda$, $\eta$.
\item \textbf{Vectors} are lowercase italic, written as columns by default: $x$, $w$, $\theta$. Coordinates are subscripted: $x_j$ is the $j$-th coordinate of $x$.
\item \textbf{Matrices} are uppercase italic: $X$, $A$, $\Sigma$. Rows index examples and columns index features in a design matrix unless stated otherwise.
\item \textbf{Random variables} are uppercase italic: $X$, $Y$, $Z$. Realizations are lowercase: $x$, $y$, $z$. Context disambiguates the dual use of $X$ as both a design matrix and a random input.
\item \textbf{Sets and spaces} use blackboard or calligraphic letters: $\mathbb{R}$, $\mathbb{R}^d$, $\mathcal{D}$ (data-generating distribution), $\mathcal{H}$ (hypothesis class).
\item \textbf{Estimates and predictions} carry a hat: $\hat{y}$, $\hat{p}$, $\hat{\theta}$, $\hat{R}_n$.
\item \textbf{Indexing.} The index $i$ ranges over examples ($i=1,\ldots,n$); the index $j$ ranges over features ($j=1,\ldots,d$); the index $k$ ranges over classes ($k=1,\ldots,K$); the index $t$ ranges over time steps.
\item \textbf{Sample versus population.} Sample quantities use sample notation ($\bar{a}$, $s_m$, $\hat{R}_n$); population or distributional quantities use $\mathbb{E}$, $\mathrm{Var}$, $R$, and named distributions. The symbol $\sigma$ is reserved for the sigmoid function and for population or noise scales; sample standard deviation is written $s_m$ to avoid the clash.
\end{itemize}

\section{Symbol Table}

\begin{table}[h]
\centering
\small
\setlength{\tabcolsep}{4pt}
\begin{tabular}{@{}L{2.0cm}L{4.7cm}L{2.6cm}@{}}
\toprule
Symbol & Meaning & Introduced in \\
\midrule
$x$ & Input feature vector for one example & Bridge~I \\
$x_j$ & The $j$-th coordinate of $x$ & Bridge~I \\
$X$ & Design matrix; rows are examples, columns are features & Bridge~I \\
$x_i$ & The $i$-th example vector; row \(i\) of \(X\) is \(x_i^\top\) & Bridge~I \\
$y$, $y_i$ & True target for an example & Bridge~II \\
$\hat{y}$, $\hat{y}_i$ & Predicted output for an example & Ch.~\ref{ch:prediction-loss-decision} \\
$\hat{p}$, $\hat{p}_i$ & Predicted probability for the positive class & Bridge~II \\
$w$ & Linear-model weight vector & Bridge~I \\
$b$ & Bias / intercept scalar & Ch.~\ref{ch:linear-models} \\
$\theta$ & Generic model parameter vector & Bridge~II \\
$f_\theta$ & Predictor function with parameters $\theta$ & Bridge~II \\
$n$ & Number of training examples & Bridges \\
$d$ & Number of features (input dimension) & Bridge~I \\
$K$ & Number of classes & Bridge~II \\
$w^\top x$ & Dot product (weighted sum) of $w$ and $x$ & Bridge~I \\
$\|x\|_1$ & $\ell_1$ norm: $\sum_j |x_j|$ & Bridge~I \\
$\|x\|_2$ & $\ell_2$ (Euclidean) norm: $\sqrt{\sum_j x_j^2}$ & Bridge~I \\
$\|x\|_\infty$ & Max-absolute-value norm: $\max_j |x_j|$ & Bridge~I \\
$\|A\|_F$ & Frobenius norm of a matrix: $\sqrt{\sum_{ij} A_{ij}^2}$ & Bridge~I \\
$\sigma(z)$ & Sigmoid function $1/(1+e^{-z})$ & Bridge~II \\
$\mathrm{softmax}(z)$ & Vector of class probabilities from scores $z$ & Bridge~II \\
$\log$ & Natural logarithm (base $e$) & Bridge~II \\
$\exp$, $e^z$ & Exponential function & Bridge~II \\
$\mathbb{P}$, $P$ & Probability of an event & Bridge~II \\
$P(Y \mid X)$ & Conditional probability of $Y$ given $X$ & Bridge~II \\
$\mathbb{E}[Z]$ & Expectation of $Z$ & Bridge~II \\
$\mathrm{Var}(Z)$ & Variance of $Z$ & Bridge~II \\
$\mathrm{Cov}(X,Y)$ & Covariance of $X$ and $Y$ & Bridge~II \\
$X \sim \mathcal{D}$ & $X$ is drawn from distribution $\mathcal{D}$ & Bridge~II \\
$\ell(y,\hat{y})$ & Pointwise loss on one example & Bridge~II \\
$R(f)$ & Population risk of predictor $f$ & Bridge~II \\
$\hat{R}_n(f)$ & Empirical risk on a sample of size $n$ & Bridge~II \\
$\mathcal{L}(\theta)$ & Training objective as a function of parameters & Bridge~II \\
$\nabla_\theta \mathcal{L}$ & Gradient of $\mathcal{L}$ with respect to $\theta$ & Bridge~II \\
$\partial \mathcal{L}/\partial \theta_j$ & Partial derivative with respect to one parameter & Bridge~II \\
$\eta$ & Learning rate / step size & Bridge~II \\
$\lambda$ & Regularization strength & Bridge~II, Ch.~\ref{ch:linear-models} \\
$\Omega(\theta)$ & Regularization penalty (e.g.\ $\|\theta\|_2^2$) & Ch.~\ref{ch:linear-models} \\
$\tau$ & Decision threshold & Bridge~II, Ch.~\ref{ch:prediction-loss-decision} \\
$\bar{a}$ & Sample mean of $a_1,\ldots,a_n$ & Bridges \\
$s_m$ & Sample standard deviation & Ch.~\ref{ch:experiments-claims} \\
$\mathrm{SE}$ & Standard error of the mean & Ch.~\ref{ch:experiments-claims} \\
TP, FP, TN, FN & Confusion-matrix counts & Ch.~\ref{ch:prediction-loss-decision} \\
$:=$, $\leftarrow$ & Definition; assignment / update & Bridges \\
\bottomrule
\end{tabular}
\end{table}

\section{Common Formulas}

The following formulas recur across the book. Each row gives the canonical form used in the text and a chapter pointer.

\begin{table}[h]
\centering
\small
\setlength{\tabcolsep}{4pt}
\begin{tabular}{@{}L{3.2cm}L{4.0cm}L{2.0cm}@{}}
\toprule
Name & Formula & Pointer \\
\midrule
Linear score & $w^\top x + b$ & Ch.~\ref{ch:linear-models} \\
Sigmoid & $\sigma(z) = 1/(1+e^{-z})$ & Bridge~II \\
Softmax & $\mathrm{softmax}(z)_k = e^{z_k}/\sum_\ell e^{z_\ell}$ & Bridge~II \\
Squared error / MSE & $\frac{1}{n}\sum_i (y_i-\hat{y}_i)^2$ & Bridge~II \\
Log loss (binary) & $-[y\log \hat{p} + (1-y)\log(1-\hat{p})]$ & Bridge~II \\
Hinge loss & $\max(0,\, 1 - y\,(w^\top x + b))$ & Ch.~\ref{ch:structure-beyond-linearity} \\
Cross-entropy (multiclass) & $-\sum_k y_k \log \hat{p}_k$ & Ch.~\ref{ch:optimization-neural-networks} \\
Empirical risk & $\hat{R}_n(f) = \frac{1}{n}\sum_i \ell(y_i, f(x_i))$ & Bridge~II \\
Regularized objective & $\hat{R}_n(f_\theta) + \lambda \Omega(\theta)$ & Ch.~\ref{ch:linear-models} \\
OLS closed form & $\hat{w} = (X^\top X)^{-1} X^\top y$ & Ch.~\ref{ch:linear-models} \\
Gradient-descent update & $\theta \leftarrow \theta - \eta\,\nabla_\theta \mathcal{L}$ & Bridge~II \\
Bayes' rule & $P(Y\mid X) = P(X\mid Y)P(Y) / P(X)$ & Bridge~II \\
Bayes-optimal threshold & $\tau = C_{FP}/(C_{FP}+C_{FN})$ & Bridge~II \\
KL divergence & $\mathrm{KL}(p\|q) = \sum_z p(z)\log\frac{p(z)}{q(z)}$ & Ch.~\ref{ch:probabilistic-models} \\
ELBO (qualitative) & $\log P(x) \geq \mathbb{E}_q[\log P(x,z)] - \mathbb{E}_q[\log q(z)]$ & Ch.~\ref{ch:generative-representations} \\
Hoeffding bound & $P(|\hat{R}_n - R| \geq \epsilon) \leq 2 e^{-2 n \epsilon^2}$ & Ch.~\ref{ch:evaluation-discipline} \\
Standard error & $\mathrm{SE} = s_m / \sqrt{S}$ & Ch.~\ref{ch:experiments-claims} \\
Boosting update (additive) & $F_{m}(x) = F_{m-1}(x) + \nu\,h_m(x)$ & Ch.~\ref{ch:structure-beyond-linearity} \\
\bottomrule
\end{tabular}
\end{table}

\section{Pointers Back to the Bridges}

For the underlying ideas behind these symbols and formulas, see the two mathematical bridges at the start of the book: Bridge~I (Linear Algebra for AI Thinking) for representation, geometry, and norms; Bridge~II (Probability, Risk, and Learning Signals) for probability, loss, empirical risk, and gradients. The bridges are the development; this appendix is the lookup. When a chapter introduces a symbol not listed here, the chapter itself is the authoritative reference.
